% VI.04.P2
\subsection*{Problem VI.04.P2 --- Products of affine algebraic sets}
\label{prob:vi04-p2}
\paragraph{Problem.}
Assume \(k\) is algebraically closed.  Let \(X\subseteq\mathbb A^m\) and
\(Y\subseteq\mathbb A^n\) be affine algebraic sets.  Prove that \(X\times Y\) is affine
algebraic and that
\[
k[X\times Y]\cong k[X]\otimes_k k[Y].
\]
If \(X\) and \(Y\) are irreducible, prove that \(X\times Y\) is irreducible.
\ifdefined\FullProblemDossiers
\paragraph{Why this problem matters.}  Products are the first place where a geometric
construction turns into a universal algebraic construction: tensor product.
\paragraph{Useful ideas before starting.}  Separate the \(x\)- and \(y\)-variables.  Use
\(I(X)k[x,y]+I(Y)k[x,y]\), then the standard quotient description of a tensor product.
Over an algebraically closed field, finitely generated reduced algebras are geometrically
reduced; for irreducible varieties their tensor product is a domain.
\paragraph{Guided hints.}  \textbf{Hint 1.} Extend the two defining ideals to
\(k[x_1,\dots,x_m,y_1,\dots,y_n]\).  \textbf{Hint 2.} A point satisfies the sum ideal iff
its two coordinates lie in \(X\) and \(Y\).  \textbf{Hint 3.} Apply
\((A/I)\otimes(B/J)\cong(A\otimes B)/(I\otimes B+A\otimes J)\).
\paragraph{Solution.}
Put \(R=k[x_1,\dots,x_m]\), \(S=k[y_1,\dots,y_n]\).  In
\(R\otimes_kS\cong k[x,y]\), the ideal
\[
K=I(X)\,k[x,y]+I(Y)\,k[x,y]
\]
has zero set exactly \(X\times Y\).  Over algebraically closed \(k\), the resulting
quotient is reduced, hence \(K=I(X\times Y)\).  Therefore
\[
\begin{aligned}
k[X\times Y]
&\cong k[x,y]/K\\
&\cong (R/I(X))\otimes_k(S/I(Y))\\
&\cong k[X]\otimes_k k[Y].
\end{aligned}
\]
If \(X,Y\) are irreducible, their coordinate rings are finitely generated domains over an
algebraically closed field and are geometrically integral; the tensor product is a domain.
Hence \(X\times Y\) is irreducible.
\paragraph{What happened mathematically.}  Independent coordinates produce a tensor
product, and irreducibility becomes absence of zero divisors.
\paragraph{Extensions.}  Identify the pullbacks of the two projections under the tensor
product description and verify the universal property of the product.
\paragraph{Provenance.}  Recovered from legacy DG4 affine Problem 9, with the
algebraically-closed hypothesis made explicit for the clean coordinate-ring statement.
\fi
